\section{K1.5 到 K2：月之暗面的技术演进}

\subsection{K1.5：强化学习技术验证}

K1.5 主要是强化学习技术路线的验证——月之暗面是较早在这条路线上投入的团队。关键发现：

\begin{importantbox}{端到端 Reward 胜过 Process Reward}
K1.5 的重要技术发现：不太需要 Process Reward Model（PRM）或 Value Function。直接使用端到端的 Reward 就可以把训练做得非常好。Process Reward 在训练过程中甚至可能有副作用。这在早期并不是非常明确的结论。
\end{importantbox}

\subsection{K2：追求更好的基础模型}

K2 的设计目标有两个核心方向：

\subsubsection{方向一：Token Efficiency——把每份数据的价值最大化}

\begin{importantbox}{为什么 Token Efficiency 比 Compute Efficiency 更重要}
高质量数据的增长非常缓慢——可以近似看作一个常数（如 30T 高质量 Token）。在这个约束下：
\begin{itemize}
    \item \textbf{Compute Efficiency}（训练更快）：虽然有价值，但不能提升智能上限——你只是更快地训练完同样的数据
    \item \textbf{Token Efficiency}（学得更好）：同样的数据，模型吸收得更多、压缩率更高、loss 降得更快——直接提升智能上限
\end{itemize}
\end{importantbox}

月之暗面通过三种方式提升 Token Efficiency：

\begin{enumerate}
    \item \textbf{Muon 优化器}：基于矩阵正交化的新优化器，取代使用了 10 年的 Adam。在 Compute Optimal 条件下，Token Efficiency 提升约 2 倍——等于学 1 份数据相当于 Adam 学 2 份
    \item \textbf{更大的稀疏度}：通过更大的稀疏度加入更多参数，参数越多，同样数据的吸收效果越好
    \item \textbf{数据 Rephrase}：对高质量数据进行改写，使同一份数据以不同形式被多次学习，提升泛化性
\end{enumerate}

\begin{knowledgebox}{Muon 优化器的技术背景}
Muon 优化器由 Keller 提出，核心思想是不把矩阵中的每个参数元素独立考虑（如 Adam 那样），而是考虑参数之间的依赖关系（dependency），通过矩阵正交化的方式实现更高效的学习。

月之暗面的贡献是将 Muon 在大规模语言模型上实际落地。早期通过 Moonlight 项目首次在一定规模的语言模型上验证，后续进一步 Scale 时发现新问题（如 MaxLogit 爆炸），通过提出新的 Clipping 方式解决。
\end{knowledgebox}

\begin{warningbox}{数据改写不等于简单复制}
单纯重复学习同一份数据会导致过拟合和泛化问题。改写的关键是引入"新的熵"——使改写后的数据在保留核心知识的同时具有一定的多样性，从而提升泛化能力。但改写方式的选择至关重要，如何引入新的信息熵仍是活跃的研究方向。
\end{warningbox}

\subsubsection{方向二：Agentic 能力与泛化}

\begin{importantbox}{Agent 泛化是当前最大挑战}
对于 Agentic 模型，当前最大的挑战是\textbf{泛化}：
\begin{itemize}
    \item 现有 RL 技术的局限：训练任务和评价指标往往是单点的（如只训练 SWE-Bench）
    \item 指标提升不代表泛化提升：SWE-Bench 分数提高，不意味着模型在其他代码任务上也好
    \item Agent 的泛化问题比对话模型更严重
    \item Benchmark 不够用或已失效
\end{itemize}
\end{importantbox}

杨植麟认为可能的解法是"用 AI 训练 AI"：让模型参与到自身的训练和对齐过程中，用更 AI-Native 的方式提升泛化性。

\subsection{K2 的研发过程}

\begin{itemize}
    \item 技术积累从一年前开始（Muon 优化器等技术需要长周期的 Scaling 实验验证）
    \item K2 本身的训练周期并不长，但前置的技术准备需要大量时间
    \item 遇到的主要挑战：Muon 优化器在大规模训练时出现 MaxLogit 爆炸问题（小规模实验中无法复现）
    \item 研究团队和训练团队是同一个团队——因为训练中的问题需要深入理解底层技术才能解决
\end{itemize}

\subsection{本章小结}

K1.5 验证了端到端 Reward 的强化学习路线。K2 聚焦两个方向：(1) Token Efficiency（Muon 优化器 + 稀疏度 + 数据改写），(2) Agentic 泛化能力。Agent 泛化是当前最大瓶颈，可能需要用 AI 训练 AI 来突破。


